
          %%%%% ~~~~~~~~~~~~~~~~~~~~ %%%%%

\section{Ratio estimators -- summary of \S5.5 and \S5.6, \cite{Sarndal1992}}

          %%%%% ~~~~~~~~~~~~~~~~~~~~ %%%%%

\subsection{Taylor linearization technique for variance estimation, \S5.5, \cite{Sarndal1992}}

          %%%%% ~~~~~~~~~~~~~~~~~~~~ %%%%%

\begin{proposition}
\mbox{}
\vskip 0.1cm
\noindent
\textbf{Setting}
\vskip 0.1cm
\noindent
Suppose:
\begin{itemize}
\item
    $p : \mathcal{P}(U) \longrightarrow [0,1]$\,
    is a sampling design satisfying:
    For each unit \,$i \in U$,\,
    its probability of selection
    \,$\pi_{i} \,:=\, \underset{s \ni i}{\sum}\;p(s)$\,
    is strictly greatly than zero.
    For each unit \,$i \in U$,\,
    its design weight is, by definition,
    \,$w_{i} \,:=\, \dfrac{1}{\pi_{i}}$.\,
\item
    For \,$k = 1, 2, \ldots, q$,\,
    \,$y^{(k)} : U \longrightarrow \Re$\,
    is an $\Re$-valued population characteristic defined on \,$U$\,
    and, we denote the population total of \,$y^{(k)}$\, by \,$Y^{(k)}$,\, i.e.,
    \,$Y^{(k)} \; := \; \underset{i \in U}{\sum}\;y^{(k)}_{i}$.\,
\item
    For \,$k = 1, 2, \ldots, q$,\,
    we denote the Horvitz-Thompson estimator of \,$y^{(k)}$\, by \,$\widehat{Y}^{(k)}$,\, i.e.,
    \begin{equation*}
    \widehat{Y}^{(k)}(s)
    \;\; := \;\;
        \underset{i \in s}{\sum}\;\dfrac{ y^{(k)}_{i} }{ \pi_{i} }\,,
    \quad\quad
    \textnormal{for each \,$s \in \mathcal{P}(U)$\, with \,$p(s) > 0$\,}
    \end{equation*}
\item
    $f : \Re^{q} \longrightarrow \Re$ is a smooth function.
\end{itemize}
\vskip 0.1cm
\noindent
\textbf{The estimation problem and its approximate solution}
\vskip 0.1cm
\noindent
\begin{itemize}
\item
    The target population parameter is:
    \begin{equation*}
    f\!\left(\,\mathbf{Y}\,\right)
    \;\; = \;\;
        f\!\left(\,Y^{(1)},\ldots,Y^{(q)}\,\right)
    \end{equation*}
\item
    We choose to use the following as the design-based estimator for \,$f\!\left(\,\mathbf{Y}\,\right)$:
    \begin{equation*}
    f\!\left(\,\widehat{\mathbf{Y}}(s)\,\right)
    \;\; = \;\;
        f\!\left(\,\widehat{Y}^{(1)}(s),\ldots,\widehat{Y}^{(q)}(s)\,\right)
    \end{equation*}
\end{itemize}
\vskip 0.1cm
\noindent
\textbf{Conclusions}
\vskip 0.1cm
\noindent
Then, the mean squared error (MSE) of the estimator
\,$f\!\left(\,\widehat{\mathbf{Y}}(s)\,\right)$\,
can be approximated as follows:
\begin{eqnarray*}
\MSE\!\left(\,
    f(\,\widehat{\mathbf{Y}}\,)
    \,\right)
& := &
    E\!\left[\,\left(\,
        f(\,\widehat{\mathbf{Y}}\,)
        \, - \,
        f(\,\mathbf{Y}\,)
        \,\right)^{2}\,\right]
\\
& {\color{red}\approx} &
    \underset{i \in U}{\sum}\;
    \underset{j \in U}{\sum}\;
    \left(\,\pi_{ij} - \pi_{i}\pi_{j}\,\right)
    \cdot
    \dfrac{u_{i}}{\pi_{i}}
    \cdot
    \dfrac{u_{j}}{\pi_{j}}
\\
& {\color{red}\approx} &
    \underset{i \in s}{\sum}\;
    \underset{j \in s}{\sum}\,
    \left(\,
        \dfrac{
            \pi_{ij} - \pi_{i}\pi_{j}
            }{
            \pi_{ij}
            }
        \,\right)
    \cdot
    \dfrac{\widehat{u}_{i}}{\pi_{i}}
    \cdot
    \dfrac{\widehat{u}_{j}}{\pi_{j}}\,,
\end{eqnarray*}
where
\begin{equation*}
u_{i}
\;\; := \;\;
    \overset{q}{\underset{k\,=\,1}{\sum}}\;
    \dfrac{\partial f}{\partial \xi_{k}}\!\left(\,
        Y^{(1)},\ldots,Y^{(q)}
        \,\right)
    \cdot
    y^{(k)}_{i}
\quad\quad
\textnormal{and}
\quad\quad
\widehat{u}_{i}
\;\; := \;\;
    \overset{q}{\underset{k\,=\,1}{\sum}}\;
    \dfrac{\partial f}{\partial \xi_{k}}\!\left(\,
        \widehat{Y}^{(1)},\ldots,\widehat{Y}^{(q)}
        \,\right)
    \cdot
    y^{(k)}_{i}
\end{equation*}
\end{proposition}
\proof
By (the multivariate version of) Taylor's Theorem, there exists a function
\,$h : \Re^{q} \longrightarrow \Re$\, such that
\begin{equation*}
\underset{\xi\,\rightarrow\mathbf{Y}}{\lim}\;h(\xi)
\; = \;
    0\,,
\quad\quad
\textnormal{and}
\quad\quad
f\!\left(\,\xi\,\right)
\; = \;
    f\!\left(\,\mathbf{Y}\,\right)
    \; + \;
    \left(\overset{{\color{white}.}}{\nabla f}\right)\!(\mathbf{Y})
    \bullet
    \left(\,\xi - \mathbf{Y}\,\right)
    \; + \;
    h(\xi)\cdot\left\Vert\,\xi - \mathbf{Y}\,\right\Vert
\end{equation*}
Hence,
\begin{eqnarray*}
f\!\left(\,\widehat{\mathbf{Y}}(s)\,\right)
& = &
    f\!\left(\,\mathbf{Y}\,\right)
    \; + \;
    \left(\overset{{\color{white}.}}{\nabla f}\right)\!(\mathbf{Y})
    \bullet
    \left(\,\widehat{\mathbf{Y}}(s) - \mathbf{Y}\,\right)
    \; + \;
    h\!\left(\,\widehat{\mathbf{Y}}(s)\,\right)
    \cdot
    \left\Vert\,
        \widehat{\mathbf{Y}}(s) - \mathbf{Y}
        \,\right\Vert
\\
& {\color{red}\approx} &
    f\!\left(\,Y^{(1)},\ldots,Y^{(q)}\,\right)
    \; + \;
    \overset{q}{\underset{k\,=\,1}{\sum}}\;
    \dfrac{\partial f}{\partial \xi_{k}}\!\left(\,
        Y^{(1)},\ldots,Y^{(q)}
        \,\right)
    \cdot
    \left(\,\widehat{Y}^{(k)}(s)\,-\,Y^{(k)}\,\right)
    % \; + \;
    % O_{p}\!\left(\,\Vert\,\widehat{\mathbf{Y}} - \mathbf{Y}\,\Vert^{2}\,\right)
\end{eqnarray*}
Hence,
\begin{eqnarray*}
\MSE\!\left(\,
    f(\,\widehat{\mathbf{Y}}\,)
    \,\right)
& := &
    E\!\left[\,\left(\,
        f(\,\widehat{\mathbf{Y}}\,)
        \, - \,
        f(\,\mathbf{Y}\,)
        \,\right)^{2}\,\right]
\\
& {\color{red}\approx} &
    E\!\left[\,\left(\,
        (\overset{{\color{white}.}}{\nabla f})(\mathbf{Y})
        \bullet
        \left(\,\widehat{\mathbf{Y}} - \mathbf{Y}\,\right)
        \,\right)^{2}\,\right]
\;\; = \;\;
    E\!\left[\,\left(\,
        (\overset{{\color{white}.}}{\nabla f})(\mathbf{Y})
        \bullet
        \widehat{\mathbf{Y}}
        \, - \,
        (\overset{{\color{white}.}}{\nabla f})(\mathbf{Y})
        \bullet
            \mathbf{Y}
        \,\right)^{2}\,\right]
\\
& = &
    \Var\!\left[\;
        \overset{{\color{white}.}}{
            (\overset{{\color{white}.}}{\nabla f})(\mathbf{Y})
            \bullet
            \widehat{\mathbf{Y}}
            }
        \;\right]
\;\; = \;\;
    \Var\!\left[\;
        \overset{{\color{white}.}}{
            \overset{q}{\underset{k\,=\,1}{\sum}}\;
            \dfrac{\partial f}{\partial \xi_{k}}\!\left(\,
                Y^{(1)},\ldots,Y^{(q)}
                \,\right)
            \cdot
            \widehat{Y}^{(k)}(s)
            }
        \;\right]
\\
& = &
    \Var\!\left[\;
        \overset{{\color{white}.}}{
            \overset{q}{\underset{k\,=\,1}{\sum}}\;
            \dfrac{\partial f}{\partial \xi_{k}}\!\left(\,
                Y^{(1)},\ldots,Y^{(q)}
                \,\right)
            \cdot
            \left(\,\underset{i \in s}{\sum}\;\dfrac{y^{(k)}_{i}}{\pi_{i}}\,\right)
            }
        \;\right]
\\
& = &
    \Var\!\left[\;
        \overset{{\color{white}.}}{
            \underset{i \in s}{\sum}\;
            \dfrac{1}{\pi_{i}}
            \left(\,
                \overset{q}{\underset{k\,=\,1}{\sum}}\;
                \dfrac{\partial f}{\partial \xi_{k}}\!\left(\,
                    Y^{(1)},\ldots,Y^{(q)}
                    \,\right)
                \cdot
                y^{(k)}_{i}
                \,\right)
            }
        \;\right]
\;\; = \;\;
    \Var\!\left[\;\,
        \overset{{\color{white}.}}{
            \underset{i \in s}{\sum}\;\dfrac{u_{i}}{\pi_{i}}
            }
        \;\right]
\\
& = &
    \underset{i \in U}{\sum}\;
    \underset{j \in U}{\sum}\;
    \left(\,\pi_{ij} - \pi_{i}\pi_{j}\,\right)
    \cdot
    \dfrac{u_{i}}{\pi_{i}}
    \cdot
    \dfrac{u_{j}}{\pi_{j}}
\\
& {\color{red}\approx} &
    \underset{i \in s}{\sum}\;
    \underset{j \in s}{\sum}\,
    \left(\,
        \dfrac{
            \pi_{ij} - \pi_{i}\pi_{j}
            }{
            \pi_{ij}
            }
        \,\right)
    \cdot
    \dfrac{\widehat{u}_{i}}{\pi_{i}}
    \cdot
    \dfrac{\widehat{u}_{j}}{\pi_{j}}
\end{eqnarray*}
\vskip -0.5cm
\qed

          %%%%% ~~~~~~~~~~~~~~~~~~~~ %%%%%

\subsection{Definition of ratio estimators}

          %%%%% ~~~~~~~~~~~~~~~~~~~~ %%%%%

\begin{definition}
\mbox{}
\vskip 0.05cm
\noindent
Suppose:
\begin{itemize}
\item
    $p : \mathcal{P}(U) \longrightarrow [0,1]$\,
    is a sampling design satisfying:
    For each unit \,$i \in U$,\,
    its probability of selection
    \,$\pi_{i} \,:=\, \underset{s \ni i}{\sum}\;p(s)$\,
    is strictly greatly than zero.
    For each unit \,$i \in U$,\,
    its design weight is, by definition,
    \,$w_{i} \,:=\, \dfrac{1}{\pi_{i}}$.\,
\item
    $y : U \longrightarrow \Re$\,
    $z : U \longrightarrow \Re$\,
    are two population characteristics and
    \,$T_{y} \,:=\, \underset{i \in U}{\sum}\;y_{i}$\,
    \,$T_{z} \,:=\, \underset{i \in U}{\sum}\;z_{i}$\,
    are their respective population totals.
% \item
%     $X^{(1)}$,\, $X^{(2)}$,\, $\ldots$\;,\, $X^{(q)}$\,
%     are a (finite) collection of ($\Re$-valued) population characteristics
%     whose respective population totals are known,
%     i.e., the following vector of population totals is known:
%     \begin{equation*}
%     T_{\mathbf{x}}
%     \;\; := \;\;
%         \left(\;\,
%             \underset{i \in U}{\sum}\,x_{i}^{(1)}
%             \,,\,
%             \underset{i \in U}{\sum}\,x_{i}^{(2)}
%             \,,\,
%             \ldots
%             \,,\,
%             \underset{i \in U}{\sum}\,x_{i}^{(q)}
%             \;\,\right)
%     \end{equation*}
% \item
%     $q_{i} > 0$,\, for each \,$i \in U$,\, has been prescribed.
\end{itemize}
The \textbf{ratio estimator} for the population ratio
\,$R_{y/z} \,:=\, \dfrac{{\color{white}.}T_{y}{\color{white}.}}{T_{z}}$\,
is defined as follows:
\begin{eqnarray*}
\widehat{R}_{y/z}(s)
& := &
    \dfrac{{\color{white}..}
        \underset{i \in s}{\sum}\;w_{i}\,y_{i}
        {\color{white}..}}{
        \underset{i \in s}{\sum}\;w_{i}\,z_{i}
        }
\end{eqnarray*}
\end{definition}

          %%%%% ~~~~~~~~~~~~~~~~~~~~ %%%%%

\vskip 0.5cm
\subsection{Explicit formulas for ratio estimators}

          %%%%% ~~~~~~~~~~~~~~~~~~~~ %%%%%

\begin{lemma}
\mbox{}
\vskip 0.01cm
\noindent
\begin{enumerate}
\item
    See last equation in \S6.4, p.234, \cite{Wolter2006} for the following formula:
    \begin{eqnarray*}
    \Var\!\left(\,\widehat{R}_{y/z}\,\right)
    & \approx &
        (\widehat{R}_{y/z})^{2}\!\left(\,
            \dfrac{\widehat{\sigma}_{\widehat{T}_{y}}}{\widehat{T}_{y}^{2}}
            \,+\,
            \dfrac{\widehat{\sigma}_{\widehat{T}_{z}}}{\widehat{T}_{z}^{2}}
            \,-\,
            2\cdot
            \dfrac{\widehat{\sigma}_{\widehat{T}_{y}\widehat{T}_{z},}}{\widehat{T}_{y}\widehat{T}_{z}}
            \,\right)
    \\
    & = &
        \left(\,
            \dfrac{{\color{white}.}
            \widehat{T}_{y}
            {\color{white}.}}{{\color{white}.}
            \widehat{T}_{z}
            {\color{white}.}}
            \,\right)^{\!2}\!\left(\,
            \dfrac{\widehat{\sigma}_{\widehat{T}_{y}}}{\widehat{T}_{y}^{2}}
            \,+\,
            \dfrac{\widehat{\sigma}_{\widehat{T}_{z}}}{\widehat{T}_{z}^{2}}
            \,-\,
            2\cdot
            \dfrac{\widehat{\sigma}_{\widehat{T}_{y}\widehat{T}_{z},}}{\widehat{T}_{y}\widehat{T}_{z}}
            \,\right)
    % \\
    % & = &
    %     \left(\,
    %         \dfrac{{\color{white}.}
    %         \widehat{T}_{y}
    %         {\color{white}.}}{{\color{white}.}
    %         \widehat{T}_{z}
    %         {\color{white}.}}
    %         \,\right)^{\!2}\!\left(\,
    %         \dfrac{\widehat{\sigma}_{\widehat{T}_{y}}}{\widehat{T}_{y}^{2}}
    %         \,+\,
    %         \dfrac{\widehat{\sigma}_{\widehat{T}_{z}}}{\widehat{T}_{z}^{2}}
    %         \,-\,
    %         2\cdot
    %         \dfrac{\widehat{\sigma}_{\widehat{T}_{y}\widehat{T}_{z},}}{\widehat{T}_{y}\widehat{T}_{z}}
    %         \,\right)
    \end{eqnarray*}
    If the sampling design is SRSWOR, the above formula simplifies to:
    \begin{eqnarray*}
    \Var\!\left(\,\widehat{R}_{y/z}\,\right)
    & \approx &
        \dfrac{1}{\widehat{T}_{z}^{2}}\cdot\left(\,
            1 - \dfrac{n}{N}
            \,\right)
        \dfrac{1}{n}
        \left(\,
            \widehat{\Var}(\widehat{T}_{y})
            \,+\,
            \widehat{R}_{y/z}^{2}\cdot\widehat{\Var}(\widehat{T}_{z})
            \,-\,
            2\cdot\widehat{R}_{y/z}\cdot\widehat{\Cov}(\widehat{T}_{y},\widehat{T}_{z})
            \,\right)
    \end{eqnarray*}
\item
    Assume the sampling design is SRSWOR.
    \begin{eqnarray*}
    \widehat{R}_{y/z}(s)
    & = &
        \dfrac{\mu_{y}}{\mu_{z}}
        \,+\,
        \dfrac{1}{\mu_{z}}\!\left(\,
            \overline{y}(s)
            \overset{{\color{white}.}}{-}
            R\cdot\overline{z}(s)
            \right)
        \,+\,
        o_{P}\!\left(\dfrac{1}{\sqrt{n}}\right),
    \quad\;\;
    \textnormal{see Equation (5.3), p.97, \cite{Wu2020}}
    \\
    & = &
        \dfrac{T_{y}}{T_{z}}
        \,+\,
        \dfrac{1}{T_{z}}\!\left(\,
            \widehat{T}_{y}(s)
            \overset{{\color{white}.}}{-}
            R\cdot\widehat{T}_{z}(s)
            \right)
        \,+\,
        o_{P}\!\left(\dfrac{1}{\sqrt{n}}\right)
    \\
    & = &
        \dfrac{T_{y}}{T_{z}}
        \,+\,
        \dfrac{1}{T_{z}}\cdot
        \underset{i \in s}{\sum}\;w_{i}\left(\,
            y_{i} - R \cdot z_{i}
            \,\right)
        \,+\,
        o_{P}\!\left(\dfrac{1}{\sqrt{n}}\right)
    \end{eqnarray*}
\end{enumerate}
\end{lemma}
\proof
\begin{enumerate}
\item
\item
\end{enumerate}
\qed

          %%%%% ~~~~~~~~~~~~~~~~~~~~ %%%%%
